\subsection{Gap Synthesis and Research Positioning}
\label{sec:gap-synthesis}

The foregoing review identifies four interconnected research gaps in laser ultrasonic denoising. First, conventional denoising methods and generic deep learning models optimize SNR or MSE objectives without explicit constraints on acoustically relevant signal features; as a result, aggressive noise suppression routinely degrades Lamb wave arrival times, amplitude ratios, and waveform morphology essential for defect detection \cite{chen2019wavelet}. Second, deep learning models trained exclusively on high-averaging data exhibit a cross-domain generalization failure when deployed on low-averaging or defect-containing specimens, where signal statistics diverge substantially from training conditions \cite{liu2022deep, zhang2020deep}. Third, the reliance on extensive signal averaging to achieve acceptable SNR imposes a linear time penalty and risks surface modification on sensitive composites, yet reduced-averaging operation remains underexplored in the literature. Fourth, while multi-domain fusion has demonstrated improved health indicator extraction in guided wave SHM \cite{perry2026health}, it remains largely underexplored for laser ultrasonic denoising; existing deep learning approaches predominantly operate in a single domain, failing to capture complementary temporal and spectral information jointly.

These gaps are individually addressed by partial solutions scattered across the literature, but no single framework concurrently tackles all four. The present work proposes PMDF-Net, a physics-constrained multi-domain fusion network that simultaneously addresses the noise suppression--feature preservation trade-off through acoustic-physics loss terms, the cross-domain generalization challenge via a defect-sensitive architecture trained across intact and flawed specimens, the reduced-averaging constraint through single-shot inference capability, and the multi-domain information gap by fusing time, frequency, and time-frequency representations within a unified framework. Section \ref{sec:methodology} describes the proposed architecture in detail.

% word count: ~230